% ------------------------------------------------------------------------------
% Slide 6: La Discusión de Fondo: Vocales Libres vs Gestos Entrenados
% ------------------------------------------------------------------------------
\begin{frame}{La Discusión de Fondo: Vocales Libres vs Gestos Entrenados}
    \begin{columns}[t]
        \begin{column}{0.48\textwidth}
            \scriptsize
            \begin{block}{Qué Miden los Sensores en la Cara}
                Los electrodos en la piel no leen el sonido, sino \textbf{los gestos que hace la cara al hablar}, que normalmente acompañan al sonido pero se pueden desacoplar.
            \end{block}
            \vspace{0.04cm}
            \begin{alertblock}{Hablar en Silencio - Sin Voz}
                Al fonoarticular sin emitir sonido:
                \begin{itemize}\setlength{\itemsep}{1pt}\setlength{\parskip}{0pt}
                    \item No hay aire saliendo ni sonido de referencia.
                    \item La persona tiende a relajar la cara o a mover la lengua adentro sin mover los labios.
                \end{itemize}
            \end{alertblock}
        \end{column}
        \begin{column}{0.48\textwidth}
            \scriptsize
            \begin{exampleblock}{Camino 1: Entrenar el Gesto Facial}
                \begin{itemize}\setlength{\itemsep}{1pt}\setlength{\parskip}{0pt}
                    \item Enseñar al usuario a marcar un gesto claro para cada vocal:
                    \begin{itemize}\setlength{\itemsep}{0pt}\setlength{\parskip}{0pt}
                        \item \textbf{/a/:} Abrir bien la boca hacia abajo.
                        \item \textbf{/i/:} Sonreír estirando hacia los costados.
                        \item \textbf{/u/:} Hacer trompita hacia adelante.
                    \end{itemize}
                    \item Funciona rápido y da clases limpias en poco tiempo.
                \end{itemize}
            \end{exampleblock}
            \vspace{0.04cm}
            \begin{block}{Camino 2: Medir la Lengua Directamente}
                \begin{itemize}\setlength{\itemsep}{1pt}\setlength{\parskip}{0pt}
                    \item Poner un sensor debajo de la mandíbula que lea la lengua.
                    \item Permite hablar con la cara relajada, pero hay que estudiar cómo cambia la fuerza al hablar en silencio.
                \end{itemize}
            \end{block}
        \end{column}
    \end{columns}
\end{frame}
